\section{Cuadricas Proyectivas}
\subsection{Cuadricas Proyectivas}

Veamos un ejemplo de una cónica proyectiva que motive este concepto.

\ejemplo{
    En $\mathbb{P}^2$ consideremos la cónica dada por la ecuación:
    \[
    C :
    x_0^2 - x_1^2 = 0 = (x_0 + x_1)(x_0 - x_1)
    \]
    Siendo así $C$ la union de 2 rectas:
    \[
    r_1 : x_0 + x_1 = 0 \quad 
    \text{y} \quad 
    r_2 : x_0 - x_1 = 0
    \]  
    En la carta afín $\{ x_0 \neq 0 \}$, con coordenadas afines $X = \frac{x_1}{x_0}$, $Y = \frac{x_2}{x_0}$, la cónica $C$ viene dada por la ecuación:
    \[
    C \cap \{ x_0 \neq 0 \} :
    1 - X^2 = 0 \implies X = \pm 1
    \]
    Que son dos rectas afines paralelas.\\
    Como antes teniamos que la intersección de $r_1$ y $r_2$ en $\mathbb{P}^2$ era el punto $(0 : 0 : 1)$, podemos tomar una carta afin que no lo contenga, por ejemplo la carta $\{ x_2 \neq 0 \}$, con coordenadas afines $X' = \frac{x_0}{x_2}$, $Y' = \frac{x_1}{x_2}$. En esta carta la cónica $C$ viene dada por la ecuación:
    \[
    C \cap \{ x_2 \neq 0 \} :
    X'^2 - Y'^2 = 0 \iff
    X' = \pm Y'
    \]
    Que son dos rectas afines que se cortan en el punto $(0,0)$.
}

En general, el polinomio (homogeneo) que define la ecuacion puede no estar determinado por las soluciones en $\mathbb{P}_{\mathbb{R}}^2$, como por ejemplo:
\[
C : x_0^2 + x_1^2 + x_2^2 = 0
\]
Cuya unica solucion en $\mathbb{P}_{\mathbb{R}}^2$ seria el punto $(0 : 0 : 0)$, que no pertenece a $\mathbb{P}_{\mathbb{R}}^2$.\\
Algo analogo le ocurre a la conica:
\[
C : x_0^2 + x_1^2 + 2x_2^2 = 0
\]
Que tampoco tiene soluciones en $\mathbb{P}_{\mathbb{R}}^2$, por lo que debemos encontrar una manera de trabajar con estas curvas que no dependa de las soluciones de sus ecuaciones.\\

\begin{definición}[Cuádrica Proyectiva]
    Una \textbf{cuádrica proyectiva} en $\mathbb{P}=\mathbb{P} (V)$, con $V$ un espacio vectorial, es una clase de equivalencia por proporcionalidad de formas bilineales simétricas:
    \[
    \hat{\varphi} : V \times V \to \mathbb{K}
    \]
    Denotaremos $\varphi$ a la cuádrica proyectiva asociada a $\hat{\varphi}$.
\end{definición}

\begin{definición}[Conjunto de ceros de una cuádrica proyectiva]
    Sea $\varphi$ una cuádrica proyectiva en $\mathbb{P}$, su conjunto de ceros es el conjunto:
    \[
    C_\varphi = 
    \mathbb{P} ( \{ \vec{u} \in V : \hat{\varphi} (\vec{u}, \vec{u}) = 0 \} )
    \]
\end{definición}

\begin{observación}
    Notese que:
    \[
    \lambda \hat{\varphi}(\mu \vec{u}, \mu \vec{u}) = 0 \quad 
    \forall \lambda, \mu \neq 0 \iff
    \hat{\varphi}(\vec{u}, \vec{u}) = 0
    \]
    Por lo que $\{ \vec{u} \in V : \hat{\varphi} (\vec{u}, \vec{u}) = 0 \}$ es una union de rectas vectoriales de $V$, y es el mismo para cualquier $\lambda \hat{\varphi}$ con $\lambda \neq 0$, es decir, depende solo de la cuádrica proyectiva $\varphi$.
\end{observación}

\subsection{Cuadricas en coordenadas}

Sea $B = \{ \vec{u_0}, \vec{u_1}, \ldots , \vec{u_n} \}$ una base de $V$, entonces la matriz de $\hat{\varphi}$ en la base $B$ es la matriz que satisface:
\[
    \hat{\varphi}(\vec{u}, \vec{v}) =
    \begin{pmatrix}
        x_0 & x_1 & \ldots & x_n
    \end{pmatrix}
    M(\hat{\varphi})_B
    \begin{pmatrix}
        y_0 \\ y_1 \\ \vdots \\ y_n
    \end{pmatrix}
\]
Siendo $\vec{u} = (x_0, x_1, \ldots , x_n)_B$ y $\vec{v} = (y_0, y_1, \ldots , y_n)_B$, por lo que:
\[
    M(\hat{\varphi})_B = (m_{ij}) \quad \text{con} \quad
    m_{ij} = \hat{\varphi}(\vec{u_i}, \vec{u_j})
\]
Así, $\hat{\varphi}$ es simetrica $\iff$ su matriz respecto de alguna base es simetrica $\iff$ su matriz respecto de cualquier base es simetrica.\\

Dada la cuadrica $\varphi$, la matriz de $\varphi$ respecto de la referencia $\mathfrak{R}$ de $\mathbb{P}$ es:
\[
    M(\varphi)_{\mathfrak{R}} :=
    M(\hat{\varphi})_B
\]
Siendo $B$ la base de $V$ asociada a la referencia $\mathfrak{R}$, y $\hat{\varphi}$ una representante de $\varphi$, por lo que $M(\varphi)_{\mathfrak{R}}$ es una matriz simetrica definida unicamente salvo proporcionalidad.\\

Así, obtenemos que el conjunto de ceros de $\varphi$ en coordenadas es:
\[
C_\varphi =
\begin{pmatrix}
    x_0 & : & x_1 & : & \ldots & : & x_n
\end{pmatrix}
\underbrace{M(\varphi)_{\mathfrak{R}}}_{M(\hat{\varphi})_B}
\begin{pmatrix}
    x_0 \\ x_1 \\ \vdots \\ x_n
\end{pmatrix}
= 0
\]

Por ejemplo, si consideramos la cuádrica proyectiva $\varphi$ en $\mathbb{P}^2(\mathbb{R})$ dada por la ecuacion:
\[
C_\varphi : x_0^2 - x_1^2 = 0 \iff
\left( x_0 : x_1 : x_2 \right)
\begin{pmatrix}
    1 & 0 & 0 \\
    0 & -1 & 0 \\
    0 & 0 & 0
\end{pmatrix}
\begin{pmatrix}
    x_0 \\ x_1 \\ x_2
\end{pmatrix}
= 0
\]

\subsection{Cuadricas en la recta proyectiva}
Cualquier forma bilineal simetrica $\hat{\varphi} : V \times V \to \mathbb{K}$, define por restricción una forma bilineal simétrica $\hat{\varphi}_{|W} : W \times W \to \mathbb{K}$, siendo $W$ un subespacio vectorial cualquiera de $V$.\\
Por tanto, una cuadrica proyectica $\varphi$ en $\mathbb{P}$ define, por restricción, cuadricas proyectivas en cualquier subespacio proyectivo $\mathbb{P}(W) \subseteq \mathbb{P}(V)$.\\

Si el subespacio es una recta proyectiva $r \subseteq \mathbb{P}(V)$, entonces con $r=\mathbb{P}(W)$, siendo $W = \langle \vec{u}_0, \vec{u}_1 \rangle$ (y por tanto $\dim(W) = 2$). Ahora tomemos la base formada por $\vec{u}_0$ y $\vec{u}_1$, que tiene asociada una referencia $\mathfrak{R}$ en $r$:
\[
    r \equiv \{[\lambda \vec{u}_0 + \mu \vec{u}_1] : 
    (\lambda : \mu) \in \mathbb{P}_\mathbb{K}^1 \} = \mathbb{P}(W)
\]

Si $\varphi$ es una cuadrica en $\mathbb{P}$, entonces $\hat{\varphi} : V \times V \to \mathbb{K}$, la ecuacion del lugar de ceros de $\varphi_{|r}$ es:
\[
0 = \hat{\varphi} (\vec{u}, \vec{u}), \quad 
\vec{u} = \lambda \vec{u}_0 + \mu \vec{u}_1 \in W \implies
\]
\[
    0 = \lambda^2 \hat{\varphi}(\vec{u}_0, \vec{u}_0) +
    2 \lambda \mu \hat{\varphi}(\vec{u}_0, \vec{u}_1) +
    \mu^2 \hat{\varphi}(\vec{u}_1, \vec{u}_1) =
    \lambda^2 a + 2 \lambda \mu b + \mu^2 c
\]

Lugar de ceros de $\varphi_{|r}$ en coordenadas respecto de la referencia $\mathfrak{R}$:

\[
C_{\varphi_{|r}} = C_\varphi \cap r : \quad 0 = a \lambda^2 + 2 b \lambda \mu + c \mu^2
\]

Discutimos las soluciones en función de $a, b, c$:

\begin{itemize}
    \item Si $a = 0$, $b = 0$ y $c = 0$, entonces $\varphi_{|W} \equiv 0$ y $C_{\varphi_{|r}} = r \iff C_\varphi$ contiene a $r$.
    \item Si $a = 0$, $b = 0$ y $c \neq 0$, entonces tenemos
    \[
        0= c \mu^2 \iff \mu = 0 \iff C_{\varphi_{|r}} = \{ (\lambda : \mu) = (1 : 0) \}
    \] 
    \item Si $a = 0$, $b \neq 0$, entonces tenemos
    \[
        0 = 2 b \lambda \mu + c \mu^2 \iff
        \mu (2 b \lambda + c \mu) = 0 \iff (\lambda : \mu) =
        (1 : 0) \text{ o } (c : -2b)
    \]
    Así $C_{\varphi_{|r}}$ tiene 2 puntos distintos.
    \item Si $a \neq 0$, entonces en este caso $\mu \neq 0$, en efecto, si $\mu = 0$ entonces $0 = a \lambda^2 \implies \lambda = 0 \implies (\lambda : \mu) = (0 : 0)$, que no es un punto de $\mathbb{P}^1_\mathbb{K}$. Por tanto, podemos dividir por $\mu^2$ y obtener:
    \[
        0 = a \left( \frac{\lambda}{\mu} \right)^2 + 2 b \left( \frac{\lambda}{\mu} \right) + c
    \]
    Tenemos una ecuación de segundo grado en $\frac{\lambda}{\mu}$, por lo que podemo tener 0, 1 o 2 soluciones en $\mathbb{K}$, dependiendo del discriminante y del cuerpo. Luego $C_{\varphi_{|r}}$ puede tener 0, 1 o 2 puntos distintos.
\end{itemize}

El lugar de cero de una cuádrica proyectiva es una recta ($\iff$ intereseción de una recta con $C_{\varphi}$) puede ser:

\begin{itemize}
    \item La recta completa (si la cuádrica contiene a la recta).
    \item Un punto (doble).
    \item Dos puntos distintos.
    \item Ningún punto (solamente si $\mathbb{K} = \mathbb{R}$).
\end{itemize}


\subsection{Polaridad asociada a una cuádrica proyectiva}

Dada una forma bilineal $\hat{\varphi} : V \times V \to \mathbb{K}$, podemos definir la aplicación lineal:

\[
    \begin{matrix}
        \hat{f} : V & \to & V^* \\
        \vec{u} & \mapsto & \hat{\varphi}(\vec{u}, \cdot)
    \end{matrix}
\]

Por ser $\hat{\varphi}$ bilineal, $\hat{f}(\vec{u}) = \hat{\varphi}(\vec{u}, \cdot)$ es un forma lineal $\implies \hat{f}(\vec{u}) \in V^*$, y además $\hat{f}$ es lineal.

Proyectivizando, si tenemos una cuádrica $\varphi$ en $\mathbb{P}(V)$, obtenemos una aplicación proyectiva:

\[
    \begin{matrix}
        f : \mathbb{P}(V) & \to & \mathbb{P}(V^*) \\
        [\vec{u}] & \mapsto & [\hat{\varphi}(\vec{u}, \cdot)]
    \end{matrix}
\]

Esta asociación está bien definida, porque $\hat{\varphi}(\lambda \vec{u}, \cdot) = \lambda \hat{\varphi}(\vec{u}, \cdot)$, por lo que definen el mismo punto en $\mathbb{P}(V^*)$.\\

\begin{definición} [Polaridad asociada a una cuádrica proyectiva]
    La aplicación proyectiva $f : \mathbb{P}(V) \dashrightarrow  \mathbb{P}(V^*)$ se denomina \textbf{polaridad asociada} a la cuádrica proyectiva $\varphi$.
\end{definición}

El centro $Z = \mathbb{P}(\ker(\hat{f})) = \mathbb{P}(\{\vec{u} \in V : \hat{\varphi}(\vec{u}, \cdot) \equiv 0 \})$ y sus puntos se donminan puntos singulares de $\varphi$.

\[
    Z = Sing(C_{\varphi}) 
\]

Obsérvese que todos los puntos singulares están en $C_{\varphi}$, pues $[\vec{u}] = A \in Z \implies \hat{\varphi}(\vec{u}, \cdot) \equiv 0 \implies \hat{\varphi}(\vec{u}, \vec{u}) = 0 \implies [\vec{u}] = A \in C_{\varphi}$.

Los puntos de $C_{\varphi} \setminus Z$ se denominan puntos regulares de $\varphi$, y así $C_{\varphi} = Reg(C_{\varphi}) \cup Sing(C_{\varphi})$.